\chapter{Sprint 1: Perception and Queue Analytics}
\chaptertoc

\clearpage
\section{Introduction}
Sprint 1 built the full path from raw video to actionable queue indicators. It started by detecting people, tracking them across frames, and filtering observations to the queue zone. Once that perception baseline was reliable, the sprint moved to queue interpretation: estimating how fast people arrive, how quickly service handles them, and how long a new arrival should expect to wait.

Perception and analytics were delivered together so that the later sprints could reuse a stable metric pipeline instead of rebuilding it from raw counts.

\section{Sprint Objectives and Backlog}
This sprint was focused on delivering two complementary layers: a reliable perception baseline for tracking people across frames within designated queue zones, and a queue analytics layer that converts tracked observations into operational metrics. The specific objectives are captured in the sprint backlog, which separates manager-facing outcomes from technical implementation concerns.

\section{Sprint Backlog}
The Sprint 1 backlog separates manager-facing user stories from technical implementation tasks. User stories represent stakeholder needs and value; technical tasks are internal infrastructure concerns. This structure ensures that traceability focuses on business outcomes while documenting all delivery work.
The user stories in Sprint 1 represent manager needs for reliable perception and analytics, as shown in Table~\ref{tab:sprint1_user_stories}.

\begin{table}[H]
    \centering
    \begingroup\setlength{\tabcolsep}{4pt}\renewcommand{\arraystretch}{0.95}\footnotesize
    \begin{tabularx}{\textwidth}{|l|X|l|X|}
        \hline
        User Story ID & User Story & Task ID & Task \\ \hline
        US1.1 & As a manager, I want to detect people in the video stream so that queue observations start from a reliable signal. & T1.1 & Configure the detector to identify people in live and recorded video. \\ \hline
        US1.2 & As a manager, I want to keep the same person identifiable across frames so that counts remain consistent. & T1.2 & Integrate tracking to preserve identity across consecutive frames. \\ \hline
        US1.3 & As a manager, I want to restrict counting to queue zones so that irrelevant movement does not affect the queue state. & T1.3 & Define queue zones as polygons and filter detections by zone. \\ \hline
        US1.6 & As a manager, I want queue rates estimated from tracked people so that the dashboard can interpret flow intensity. & T1.6 & Group tracked events into short windows and compute arrival and service rates. \\ \hline
        US1.7 & As a manager, I want wait time computed adaptively so that the result remains useful in real conditions. & T1.7 & Apply the M/M/1 formula when arrivals are below service capacity and fallback to occupancy-based estimation otherwise. \\ \hline
        US1.9 & As a manager, I want a clear metric output so that the dashboard can reuse these values without reprocessing them. & T1.9 & Publish rate estimates, wait time, and people-in-zone occupancy. \\ \hline
    \end{tabularx}
    \caption{Sprint 1 user stories}
    \label{tab:sprint1_user_stories}
    \endgroup
\end{table}

Supporting technical tasks address internal implementation requirements that enable these user stories, as detailed in Table~\ref{tab:sprint1_technical_tasks}.

\begin{table}[H]
    \centering
    \begingroup\setlength{\tabcolsep}{4pt}\renewcommand{\arraystretch}{0.95}\footnotesize
    \begin{tabularx}{\textwidth}{|l|X|l|X|}
        \hline
        Task ID & Task & Subtask ID & Subtask \\ \hline
        T1.4 & Connect capture, detection, tracking, zone filtering, and display into one continuous flow. & ST1.4.1 & Implement runtime orchestration that integrates all perception components. \\ \hline
        T1.5 & Expose runtime settings so that the pipeline can be adapted without code changes. & ST1.5.1 & Provide configurable parameters for frame handling and visualization. \\ \hline
        T1.8 & Add guards for empty windows, low counts, and bounded fallback output to protect calculations from division-by-zero and noise. & ST1.8.1 & Implement defensive bounds checking in the analytics layer. \\ \hline
    \end{tabularx}
    \caption{Sprint 1 technical tasks}
    \label{tab:sprint1_technical_tasks}
    \endgroup
\end{table}

\section{Design}

\subsection{Perception and Analytics Flow}
The Sprint 1 workflow from frame ingestion to queue metrics output occurs at two levels: the overall perception pipeline and the detailed queue analytics transformation. The design must address several runtime conditions and operational scenarios that arise during real-world queue monitoring.

At the perception level, the system handles three key scenarios: when a person enters the queue zone, they are treated as a possible arrival; when the same person appears in consecutive frames, the system maintains stable identification to avoid double-counting; and when people appear outside the zone, they are filtered out to keep measurements focused on the queue area.

At the queue-operation level, the system must respond gracefully to three distinct conditions. In a normal queue, arrivals remain below service capacity, allowing the M/M/1 formula to provide precise rate-based estimates. When arrivals approach service capacity, the system seamlessly switches to a conservative occupancy-based estimate. When no service events are available, the output remains bounded and safe for downstream processing. This adaptive logic ensures that a single metric responds correctly to real-world conditions without requiring separate state signals.

At the high level, Figure~\ref{fig:sprint1_activity} shows the end-to-end flow from a video source (local file or live camera stream) through detection, tracking, zone filtering, and finally to queue analysis that computes arrival rate, service rate, and expected wait time.
\begin{figure}[!htbp]
    \centering
    \includegraphics[width=0.60\textwidth]{\detokenize{chapter3_activity.pdf}}
    \caption{Sprint 1 activity flow from frame ingestion to queue metrics output.}
    \label{fig:sprint1_activity}
\end{figure}

At the detailed level, Figure~\ref{fig:sprint1_analytics_activity} focuses on the queue analytics transformation that converts tracked people inside the zone into the metric payload consumed by the dashboard and later automation layers. The analytics step compares current tracked identifiers with the previous observation window to identify arrivals and departures. These events are stored in short windows, then converted into rates ($\lambda$ and $\mu$). Before applying the wait-time formula, the system checks whether enough events were observed, whether $\mu$ is positive, and whether arrivals remain below service capacity. If these guards are satisfied, it uses the M/M/1 expression; otherwise, it falls back to a bounded occupancy-based estimate.
\begin{figure}[!htbp]
    \centering
    \includegraphics[width=0.72\textwidth]{\detokenize{chapter3_activity_Analytics.pdf}}
    \caption{Sprint 1 analytics activity diagram for queue metric generation.}
    \label{fig:sprint1_analytics_activity}
\end{figure}

\section{Implementation}

\subsection{Perception Pipeline}
The perception pipeline integrates three core components: detection, tracking, and zone filtering. The detection layer identifies people in each frame using YOLO and is tuned to balance speed and accuracy, ensuring the system remains responsive without losing critical observations. The detection output is then fed to the tracking layer, which preserves person identity across consecutive frames using ByteTrack. This identity continuity is essential for reliable arrival and departure counting, as it prevents double-counting from single-frame detections alone. Finally, zone-aware filtering restricts all measurements to a polygon-defined queue area, eliminating unrelated movement elsewhere in the scene and ensuring the metric reflects only the relevant part of the frame.

The pipeline integration is orchestrated in the runtime layer, which supports handling diverse input sources (local files, webcams, and live camera streams), adjusting frame-processing frequency to maintain system responsiveness, rendering overlays to display detections, tracker identities, queue zones, and computed metrics, and operating without visual output when needed for headless integration with downstream systems.

\subsection{Queue Event Windows and Rate Estimation}
The queue-analysis step uses short time windows to count arrivals and departures. Let $N_a$ and $N_d$ be the number of events observed in windows $T_a$ and $T_d$:
\[
\lambda = \frac{N_a}{T_a}, \qquad \mu = \frac{N_d}{T_d}
\]
From these counts, the system estimates the arrival rate and the service rate. To avoid unstable values at the beginning of a run, the implementation waits until enough events have been observed before returning non-zero rates.

\subsection{Wait-Time Logic and Adaptive Estimation}
The system estimates wait time using an adaptive formula that responds to queue conditions:
\[
W = \begin{cases}
\dfrac{1}{\mu - \lambda} & \text{if } \lambda < \mu \text{ and } \mu > 0 \\
\dfrac{L}{\mu} & \text{otherwise}
\end{cases}
\]
where $L$ is the current number of people in the zone. We use the M/M/1 sojourn-time expression because the project models each monitored lane as a single service channel with one effective queue, which matches the cashier-style setting studied in the report. Other models were less appropriate here: M/M/c would assume multiple parallel servers, which is not the default lane setup, and Little's Law is useful for validation but does not directly produce a forward wait-time estimate from the instantaneous rates alone.

The adaptive formula ensures the output always remains safe and interpretable: when arrivals are below service capacity, the formula is precise; when arrivals exceed capacity, it gracefully degrades to a conservative estimate based on current queue size. When no service events are observed, the implementation returns zero.

\subsection{Runtime Metric Contract}
The final metric output includes:
\begin{itemize}
    \item point estimates ($\lambda$, $\mu$, $W$),
    \item people-in-zone occupancy.
\end{itemize}
These values are then reused by logs, overlays, webhook payloads, and dashboard events \cite{opencv_docs,gross_queueing}.

\section{Conclusion}
Sprint 1 delivered the complete path from raw video to operational queue indicators, establishing a trustworthy metric layer for downstream systems. The integration of detection, tracking, and zone filtering produced a reliable perception baseline, while the queue analytics layer converted tracked observations into arrival rates, service rates, and adaptive wait-time estimates.

The key technical achievement was the adaptive wait-time formula, which remains interpretable and safe across all operating conditions: it provides precise estimates when arrivals are below service capacity, gracefully degrades to occupancy-based estimates when the queue grows, and returns bounded values when service events are unavailable. This design ensures that downstream layers—the dashboard, persistence, and automation—can operate on structured and stable signals rather than raw detection counts.

By connecting perception to interpretation in a single delivery, the sprint reduced false continuity problems through tracking and zone logic, while the queue model provided the necessary foundation for later integration with web monitoring and automated alerting workflows. The metric contract established in this sprint is designed to be reused consistently throughout the system's remaining layers.


